%=============================================================
%  Chapter 6 --- Iteration 2 : AI & Personalisation
%=============================================================
\chapter{Implementation --- Iteration 2: Artificial Intelligence and Personalisation}

\section*{Introduction}
The second iteration enriches the platform with an artificial intelligence layer
composed of two complementary modules: a \textbf{RAG chatbot} capable of answering
customer questions about products, and a \textbf{recommendation engine} driven by
behavioral signal tracking.
An onboarding questionnaire completes the feature set by personalizing the
experience from the very first visit.

%-------------------------------------------------------------
\section{Sprint Planning}
%-------------------------------------------------------------

\subsection{Sprint Goal}

\begin{tcolorbox}[colback=blue!5, colframe=blue!40, title=Sprint Goal --- Iteration 2]
Give every customer a personalized experience through a conversational assistant
grounded in the real catalogue and recommendations tailored to their skin profile
and browsing behavior.
\end{tcolorbox}

\subsection{Sprint Backlog}

\begin{table}[H]
\centering
\caption{User stories --- Iteration 2}
\begin{tabularx}{\textwidth}{|c|X|c|c|}
\hline
\rowcolor{gray!20}
\textbf{ID} & \textbf{User Story} & \textbf{Priority} & \textbf{Points} \\
\hline
US-11 & As a customer, I want to answer a profile questionnaire (skin type, goals) on first login. & Must & 5 \\
\hline
US-12 & As a customer, I want to ask a chatbot questions about products and get accurate answers based on the catalogue. & Must & 13 \\
\hline
US-13 & As a customer, I want the chatbot to refuse manipulation attempts (prompt injection). & Must & 8 \\
\hline
US-14 & As a customer, I want to receive product recommendations adapted to my skin type and browsing habits. & Must & 8 \\
\hline
US-15 & As a customer, I want to add products to a wishlist and view it later. & Should & 3 \\
\hline
US-16 & As a customer, I want to leave a review (rating + comment) on a product I have purchased. & Should & 5 \\
\hline
\end{tabularx}
\end{table}

%-------------------------------------------------------------
\section{Design}
%-------------------------------------------------------------

\subsection{RAG Service Architecture}

The AI service is decoupled from the main backend as a standalone
Python/FastAPI microservice, communicating with the Node.js server over
internal HTTP. This separation allows the AI stack (model, vector store) to
evolve independently of the business logic.

\begin{figure}[H]
  \centering
  \includegraphics[width=\textwidth]{figures/arch_component.png}
  \caption{Component architecture including the RAG service (see Chapter 3)}
  \label{fig:arch_rag_ref}
\end{figure}

\subsection{RAG Pipeline}

The RAG pipeline runs in four stages:

\begin{enumerate}
  \item \textbf{Indexing} --- product sheets (name, description, ingredients,
        skin type) are transformed into vectors by the
        \texttt{all-MiniLM-L6-v2} model and stored in the vector store.
  \item \textbf{Retrieval} --- the user's question is encoded, and the
        \textit{k} most semantically similar documents are retrieved.
  \item \textbf{Augmentation} --- the retrieved documents form the context
        injected into the system prompt sent to the LLM.
  \item \textbf{Generation} --- Ollama (\texttt{llama3}) produces a response
        grounded in the documents, reducing hallucinations.
\end{enumerate}

\subsection{Sequence Diagrams}

The chatbot and recommendation engine interaction flows are illustrated by
the \textit{seq\_chatbot} and \textit{seq\_recommendations} diagrams
presented in Chapter 3.

%-------------------------------------------------------------
\section{Implementation}
%-------------------------------------------------------------

\subsection{Onboarding Questionnaire}

On first login, a modal prompts the customer to select their skin type
(normal, dry, oily, combination, sensitive) and cosmetic goals
(hydration, anti-aging, radiance, etc.).
These preferences are persisted in the user profile and used as a
priority filter in the recommendation engine.

\begin{lstlisting}[language=Java, caption=Saving the onboarding profile (OnboardingModal.tsx)]
const handleComplete = async (profile: OnboardingProfile) => {
  try {
    await usersApi.updateProfile(profile);
    onComplete();   // close modal only on success
  } catch {
    setSaveError('Something went wrong. Please try again.');
  }
};
\end{lstlisting}

\begin{figure}[H]
  \centering
  \includegraphics[width=0.8\textwidth]{figures/screen_onboarding.png}
  \caption{Onboarding modal --- skin profile questionnaire}
  \label{fig:screen_onboarding}
\end{figure}

%-------------------------------------------------------------
\subsection{RAG Chatbot}
%-------------------------------------------------------------

\subsubsection{Prompt Injection Guard}

Before being forwarded to the RAG service, every request passes through
three filtering levels implemented in \texttt{RagHttpsClient}:

\begin{enumerate}
  \item \textbf{Length check} --- any message exceeding 500 characters is
        immediately rejected (\texttt{400 Bad Request}).
  \item \textbf{Regex patterns} --- a set of regular expressions detects classic
        injection attempts (\textit{ignore previous instructions},
        \textit{jailbreak}, \textit{system prompt leak}, etc.).
  \item \textbf{ML classifier} --- a lightweight classifier assigns a confidence
        score; above the threshold, the request is blocked.
\end{enumerate}

\begin{lstlisting}[language=Java, caption=Triple guard in RagHttpsClient.ts]
if (message.length > 500)
  return { answer: null, blocked: true, reason: 'too_long' };

if (INJECTION_PATTERNS.some(p => p.test(message)))
  return { answer: null, blocked: true, reason: 'regex_match' };

const score = await this.classifyInjection(message);
if (score > INJECTION_THRESHOLD)
  return { answer: null, blocked: true, reason: 'ml_classifier' };
\end{lstlisting}

\subsubsection{RAG Service (Python / FastAPI)}

The service exposes a single \texttt{POST /query} endpoint that orchestrates
the vector search and the LLM call.

\begin{lstlisting}[language=Python, caption=RAG pipeline --- main.py (FastAPI)]
@app.post("/query")
async def query(req: QueryRequest):
    # 1. Encode the question
    embedding = embedder.encode(req.question)
    # 2. Vector search (top-5)
    docs = vector_store.search(embedding, k=5)
    context = "\n\n".join(d.content for d in docs)
    # 3. Call Ollama with context
    response = ollama.chat(model="llama3", messages=[
        {"role": "system", "content": SYSTEM_PROMPT + context},
        {"role": "user",   "content": req.question}
    ])
    return {"answer": response["message"]["content"]}
\end{lstlisting}

\begin{figure}[H]
  \centering
  \includegraphics[width=\textwidth]{figures/screen_chatbot.png}
  \caption{Chatbot widget --- answer grounded in the product catalogue}
  \label{fig:screen_chatbot}
\end{figure}

\subsubsection{Client Integration}

The chatbot is accessible through a floating widget available on all pages.
Input is disabled while the response is loading, and an explicit message
is shown when a request is detected as injection.

\begin{figure}[H]
  \centering
  \includegraphics[width=0.65\textwidth]{figures/screen_chatbot_blocked.png}
  \caption{Blocked message displayed when injection is detected}
  \label{fig:screen_chatbot_blocked}
\end{figure}

%-------------------------------------------------------------
\subsection{Recommendation Engine}
%-------------------------------------------------------------

\subsubsection{Behavioral Signal Tracking}

Every significant interaction by an authenticated customer generates a signal
persisted in the \texttt{behaviour\_signals} table:

\begin{table}[H]
\centering
\caption{Behavioral signal types and their weights}
\begin{tabular}{|l|l|c|}
\hline
\rowcolor{gray!20}
\textbf{Signal} & \textbf{Trigger} & \textbf{Weight} \\
\hline
\texttt{VIEW}       & Product page visited      & 1 \\
\hline
\texttt{WISHLIST}   & Added to wishlist         & 3 \\
\hline
\texttt{CART\_ADD}  & Added to cart             & 5 \\
\hline
\texttt{PURCHASE}   & Order confirmed           & 10 \\
\hline
\end{tabular}
\end{table}

\subsubsection{Recommendation Scoring}

Recommendations are computed server-side on every call to
\texttt{GET /api/products?recommended=true}.
PostgreSQL aggregates signals per product for the current user and sorts
results by descending affinity score.
The onboarding skin type profile is applied as a priority filter.

\begin{lstlisting}[language=SQL, caption=Simplified recommendation query]
SELECT p.*, COALESCE(SUM(bs.weight), 0) AS affinity_score
FROM products p
LEFT JOIN behaviour_signals bs
       ON bs.product_id = p.id AND bs.user_id = $userId
WHERE p.is_active = true
  AND p.skin_type @> $skinTypeFilter::jsonb
GROUP BY p.id
ORDER BY affinity_score DESC, p.created_at DESC
LIMIT 20;
\end{lstlisting}

\begin{figure}[H]
  \centering
  \includegraphics[width=\textwidth]{figures/screen_recommendations.png}
  \caption{Personalized recommendations section on the homepage}
  \label{fig:screen_recommendations}
\end{figure}

%-------------------------------------------------------------
\subsection{Wishlist}
%-------------------------------------------------------------

The wishlist is persisted in the database and synchronized with the client
Zustand store. A product can be added from the product page or from the
catalogue grid via the heart button, accessible to logged-in users only.

\begin{figure}[H]
  \centering
  \includegraphics[width=\textwidth]{figures/screen_wishlist.png}
  \caption{Wishlist page}
  \label{fig:screen_wishlist}
\end{figure}

%-------------------------------------------------------------
\subsection{Customer Reviews}
%-------------------------------------------------------------

A customer can submit a review (1--5 star rating + comment) only on a product
present in their past orders, ensuring review authenticity.
The verification is performed server-side in \texttt{ReviewController}.

\begin{figure}[H]
  \centering
  \includegraphics[width=0.85\textwidth]{figures/screen_reviews.png}
  \caption{Customer reviews section on the product page}
  \label{fig:screen_reviews}
\end{figure}

%-------------------------------------------------------------
\section{Sprint Review}
%-------------------------------------------------------------

\subsection{Delivered Features}

\begin{table}[H]
\centering
\caption{Iteration 2 results}
\begin{tabular}{|c|l|c|}
\hline
\rowcolor{gray!20}
\textbf{ID} & \textbf{User Story} & \textbf{Status} \\
\hline
US-11 & Onboarding questionnaire & \textcolor{green!60!black}{\checkmark\ Done} \\
\hline
US-12 & RAG chatbot & \textcolor{green!60!black}{\checkmark\ Done} \\
\hline
US-13 & Prompt injection guard & \textcolor{green!60!black}{\checkmark\ Done} \\
\hline
US-14 & Personalized recommendations & \textcolor{green!60!black}{\checkmark\ Done} \\
\hline
US-15 & Wishlist & \textcolor{green!60!black}{\checkmark\ Done} \\
\hline
US-16 & Customer reviews & \textcolor{green!60!black}{\checkmark\ Done} \\
\hline
\end{tabular}
\end{table}

\subsection{Challenges}

\begin{itemize}
  \item \textbf{Local LLM latency}: Ollama CPU inference introduces 3--8 seconds
        per request. An animated loading indicator and streaming response
        (Server-Sent Events) were added to improve perceived performance.
  \item \textbf{Indexing quality}: early pipeline versions returned irrelevant
        documents due to poor chunking. Splitting product sheets by section
        (ingredients, benefits, usage) significantly improved answer accuracy.
  \item \textbf{Classifier false positives}: legitimate questions containing
        technical terms (\textit{ignore}, \textit{system}) were being blocked.
        The classifier threshold was calibrated against a test set representative
        of cosmetics vocabulary.
\end{itemize}

\section*{Conclusion}
Iteration 2 transformed \textbf{LUMINA} from a standard online store into an
intelligent platform.
The RAG chatbot provides conversational access to the catalogue with robust
protection against abuse, while the behavioral signal-driven recommendation
engine personalizes the experience without relying on an external machine
learning model.
Iteration 3 will finalize the platform by strengthening administration tools
and security mechanisms.
